\chapter{Türkçe--İngilizce Terim Sözlüğü}\label{ek:sozluk}

Mantık ve tip kuramı terimlerinin Türkçe karşılıkları henüz tam olarak
yerleşmiş değildir. Bu kitapta aşağıdaki karşılıklar tutarlı biçimde
kullanılmaktadır. Parantez içindeki karşılıklar da literatürde rastlanan
seçeneklerdir.

\begin{center}
\renewcommand{\arraystretch}{1.25}
\rowcolors{2}{tablosatir}{white}
\begin{tabularx}{\linewidth}{>{\raggedright\arraybackslash}X >{\raggedright\arraybackslash\itshape}X}
\rowcolor{tabloust}
\tablobaslik{Türkçe} & \tablobaslik{İngilizce} \\
adla / değerle çağırma & call-by-name / call-by-value \\
$\alpha$-denklik & alpha-equivalence \\
alt formül özelliği & subformula property \\
alt küme & subset \\
alt tip & subtype \\
alt yapısal mantık & substructural logic \\
ana formül & principal formula \\
anlambilim & semantics \\
ara değer (interpolasyon) & interpolation \\
ardıl & successor \\
ardışık hesap & sequent calculus \\
ardışık (sekant) & sequent \\
bağımlı çift & dependent pair \\
bağımlı fonksiyon & dependent function \\
bağımlı tip & dependent type \\
bağımlı tip kuramı & dependent type theory \\
bağlı değişken & bound variable \\
$\beta$-indirgeme & beta-reduction \\
bileşke & composition \\
birebir & injective \\
birleşebilirlik & confluence, Church--Rosser property \\
birleşim & union \\
bölüm tipi & quotient type \\
büyük eleme & large elimination \\
çelişki ile ispat & proof by contradiction, reductio ad absurdum \\
çelişmezlik ilkesi & law of non-contradiction \\
çıkarım kuralı & inference rule \\
çift değilleme çevirisi & double-negation translation \\
çift değilleme eleme & double negation elimination \\
çok sonuçlu ardışık & multiple-conclusion sequent \\
çözünürlük & resolution \\
daraltma & contraction \\
de Bruijn indisi & de Bruijn index \\
değiştirme & exchange \\
denklik & equivalence, biconditional \\
denklik ilişkisi & equivalence relation \\
denklik sınıfı & equivalence class \\
derin gömme & deep embedding \\
dışsallık & extensionality \\
doğal çıkarım & natural deduction \\
doğruluk değerlemesi & valuation \\
doğruluk tablosu & truth table \\
doğrusal mantık & linear logic \\
dolambaç & detour, maximal formula \\
durum ayrımı & proof by cases \\
elaboratör & elaborator \\
eleme kuralı & elimination rule \\
eşleme & bijection \\
evren & universe \\
evren hiyerarşisi & universe hierarchy \\
evrensel niceleyici & universal quantifier \\
ex falso (çelişkiden her şey) & ex falso quodlibet, principle of explosion \\
fonksiyon dışsallığı & function extensionality \\
geçerlilik (sağlamlık) & soundness \\
geri dönme & backtracking \\
geriye / ileriye akıl yürütme & backward / forward reasoning \\
giriş kuralı & introduction rule \\
güçlü normalleşme & strong normalization \\
güçlü tümevarım & strong induction \\
güvenilir çekirdek & trusted kernel \\
hedef & goal \\
hesaplanamaz & noncomputable \\
hipotetik kıyas & hypothetical syllogism \\
impredikatif & impredicative \\
indeks & index \\
İnşalar Hesabı & Calculus of Constructions \\
ispat arama & proof search \\
ispat asistanı & proof assistant \\
ispat karmaşıklığı & proof complexity \\
işaretli atom & signed atom, literal \\
iyi sıralama ilkesi & well-ordering principle \\
iyi tanımlılık & well-definedness \\
kabul edilebilir kural & admissible rule \\
kanıt-ilgisizliği & proof irrelevance \\
kanonik değerleme & canonical valuation \\
karar verilebilir & decidable \\
karşı model & countermodel \\
kaydırma & shifting \\
kesişim & intersection \\
kesme eleme & cut elimination \\
kesme kuralı & cut rule \\
kısmi sıralama & partial order \\
kip mantığı & modal logic \\
kombinatör & combinator \\
kontrapozisyon & contraposition \\
köşegen yöntemi & diagonal argument \\
kural tümevarımı & rule induction \\
kurucu & constructor \\
lif & fibre \\
Lindenbaum önsavı & Lindenbaum's lemma \\
maksimal tutarlı küme & maximal consistent set \\
mantıksal sonuç & logical consequence \\
metamantık & metalogic \\
model varlığı & model existence \\
motive & motive \\
nesne dili & object language \\
niceleyici & quantifier \\
normal türetim & normal derivation \\
normalleştirme & normalization \\
önerme & proposition \\
önerme dışsallığı & propositional extensionality \\
önermesel eşitlik & propositional equality \\
örten & surjective \\
öz değişken koşulu & eigenvariable condition \\
özdeşlik tipi & identity type \\
özyineleme & recursion \\
özyineleyici & recursor \\
parametre & parameter \\
pozitiflik koşulu & positivity condition \\
predikatif & predicative \\
redeks & redex \\
sağ / sol kural & right / left rule \\
sağlanabilir & satisfiable \\
seçim aksiyomu & axiom of choice \\
serbest değişken & free variable \\
sezgicilik & intuitionism \\
sığ gömme & shallow embedding \\
soyutlama & abstraction \\
sözdizim ağacı & syntax tree \\
sözdizimi & syntax \\
taktik & tactic \\
tamlık & completeness \\
tanık & witness \\
tanık sabiti & witness constant, Henkin constant \\
tanım evreni & domain of discourse \\
tanımsal eşitlik & definitional equality \\
taşıma & transport \\
tekdeğerlilik & univalence \\
tekdüzelik & monotonicity \\
temsilci & representative \\
terim modeli & term model \\
ters matematik & reverse mathematics \\
tersinir kural & invertible rule \\
tıkızlık & compactness \\
tip ailesi & type family \\
tip denetleyici & type checker \\
tip korunumu & subject reduction, type preservation \\
totoloji & tautology \\
tutarlılık & consistency \\
tümevarım & induction \\
tümevarım hipotezi & induction hypothesis \\
tümevarımsal aile & inductive family \\
tümevarımsal tip & inductive type \\
tümevarımsal yüklem & inductive predicate \\
tümleyen & complement \\
türetilebilirlik & derivability \\
türetilmiş kural & derived rule \\
türetim & derivation \\
uygulama & application \\
üçüncü hâlin olanaksızlığı & law of excluded middle \\
üst dil & metalanguage \\
varoluşsal niceleyici & existential quantifier \\
yakalama & variable capture \\
yansımalı-geçişli kapanış & reflexive-transitive closure \\
yapıcı (sezgici) mantık & constructive (intuitionistic) logic \\
yapısal kural & structural rule \\
yapısal özyineleme & structural recursion \\
yapısal tümevarım & structural induction \\
yerellik & locality, coincidence lemma \\
yerine koyma & substitution \\
yoğunlamalı & intensional \\
yüklem & predicate \\
zayıf / güçlü tamlık & weak / strong completeness \\
zayıflatma & weakening \\
\end{tabularx}
\end{center}
